% =========================================================================
% A3 — Additional ablations.
%
% Fragment for "Frozen Basis Time Series Forecasting" (camera-ready rewrite
% of NeurIPS submission 30981). The assembler inserts \appendix.
%
% Governing documents: paper_rewrite_plan.md §6 App D subplan (active_g
% ablation = the dagger evidence; tiered-offset evidence; boundary test;
% rectifier diagnostic; divergence-mechanism detail; nonlinearity-placement
% study), §11 corrections ledger item 5 ("tiering helps" does not survive),
% §9 register (boundary test reported). Study-local encoded strings are
% replaced by descriptive labels per notation spec §6b. No sum_losses or
% cross-normalization comparison appears.
% =========================================================================

\section{Additional Ablations}
\label{app:ablations}

This appendix collects the matched-seed evidence for every empirical rule
cited in Sections~\ref{sec:results:stability}
and~\ref{sec:setup:protocols} and in Appendix~\ref{app:hyperparams:block}.
All numbers come verbatim from the analysis reports released with the
per-run CSVs, and the strict health filter of
Section~\ref{sec:setup:stats} is applied throughout.

\subsection{The forecast-head activation: definition, instability, and
per-period effect}
\label{app:ablations:activeg}

\paragraph{Definition.} The standard generic block produces its backcast
and forecast through a pair of unconstrained linear projections from the
shared trunk state,
\[
  \hat{\mathbf{b}} = \mathbf{W}_b \mathbf{h},
  \qquad
  \hat{\mathbf{f}} = \mathbf{W}_f \mathbf{h},
\]
with $\mathbf{h} \in \mathbb{R}^{U}$ the trunk output. The extension swept
in this paper applies a nonlinear post-activation to these projections
before they enter the frozen basis,
\[
  \hat{\mathbf{b}} = \phi(\mathbf{W}_b \mathbf{h}),
  \qquad
  \hat{\mathbf{f}} = \phi(\mathbf{W}_f \mathbf{h}),
\]
where $\phi$ is the trunk activation (ReLU by default). The
$\dagger$ variant activates only the forecast projection, leaving the
backcast head linear; activating both heads is statistically
indistinguishable at $n = 200$ on M4-Yearly ($\Delta = -0.023$, $p = 0.154$)
and Tourism-Yearly ($\Delta = +0.066$, $p = 0.087$), so the forecast-head-only
form is preferred because it isolates the intervention. The parameter
defaults to off, preserving the linear output of
\citet{oreshkin2020nbeats}; paper-faithful claims exclude $\dagger$ cells.

\paragraph{Observed instabilities.} In a large-scale convergence study
($n = 200$ seeds per configuration, M4-Yearly), overparameterised plain
generic stacks --- most visibly the thirty-stack baseline --- showed a
pronounced bimodal failure mode: a substantial fraction of random seeds
converged to a degenerate fixed point with sMAPE an order of magnitude
above the well-converged subset. The same pathology appeared at extreme
rates on the small-data Milk series, where 63\% of six-stack seeds and
64\% of ten-stack seeds were classified as stuck (sMAPE above 5), and at
4.5\% on M4-Yearly. Setting the activation on either head eliminated the
stuck-mode failure \emph{entirely} in every configuration tested: the Milk
rates fell to 0\%, as did the M4-Yearly rate
(Table~\ref{tab:active_g_ablation}). The convergence-reliability gain is
unambiguous across every dataset examined.

\begin{table}[h]
\centering
\small
\caption{Forecast-head activation effect by dataset, two-sided
Mann--Whitney $U$. ``Quality cost'' is the absolute sMAPE delta on the
converged subset of seeds; the last column is the stuck-seed failure-rate
reduction. ``Converged subset'' applies a dataset-specific stuck-mode
filter beyond the strict health filter of Section~\ref{sec:setup:stats}
(sMAPE $< 20$ on M4-Yearly, $< 30$ on Tourism-Yearly, $< 5$ on Milk) so
that bimodal stuck seeds do not contaminate the converged-mean estimate.}
\label{tab:active_g_ablation}
\begin{tabular}{@{}lllrrl@{}}
\toprule
Dataset & Baseline & With activation & $\Delta$ & MWU $p$ & Convergence effect \\
\midrule
Milk, six stacks & $1.685$ ($n{=}37$) & $2.500$ ($n{=}100$) & $+0.815$ & $<10^{-6}$ & $63\% \to 0\%$ \\
Milk, ten stacks & $1.797$ ($n{=}36$) & $2.455$ ($n{=}100$) & $+0.659$ & $<10^{-5}$ & $64\% \to 0\%$ \\
M4-Yearly        & $13.616$ ($n{=}191$) & $13.669$ ($n{=}200$) & $+0.053$ & $<10^{-4}$ & $4.5\% \to 0\%$ \\
Tourism-Yearly   & $21.539$ ($n{=}197$) & $21.374$ ($n{=}200$) & $-0.165$ & $1.8\times10^{-4}$ & $1.5\% \to 0\%$ \\
\bottomrule
\end{tabular}
\end{table}

\paragraph{Per-period verdict on M4 paper-sample.} Hourly: 9 of 9 matched
pairs improve under the activation (5 significant at $p < 0.05$); the
canonical leaderboard delta is $-0.148$ sMAPE, the thirty-stack
interpretable baseline scoring 8.758 with the activation against 8.906
without. Yearly, Quarterly, Monthly, Weekly, and Daily: the activation ties
or loses on every schedule cell tested. The recommendation is a targeted
intervention, not a global default: keep it off for paper-faithful
reproduction, switch it on for M4-Hourly, and consider it elsewhere only
when bimodal convergence failures dominate and convergence reliability
outranks converged-subset accuracy.

\paragraph{A working hypothesis: unactivated linear outputs as a source of
instability.} That a single activation after the output projections
eliminates convergence failures with no other architectural change invites
a mechanistic reading, which we offer as a working hypothesis rather than a
theoretical claim. In the doubly residual stack, an unbounded linear final
projection places no inherent limit on early-training backcast or forecast
magnitudes. A large-magnitude backcast at initialisation over-subtracts the
residual; the corrupted stream is then amplified by the next block's own
unconstrained projections, and in deep stacks the cascade can compound past
recovery by standard Adam under any fixed schedule. An activation bounds
the projected signal --- ReLU restricts it to the non-negative half-space
--- and acts as a soft amplitude regulariser throughout early training. The
benefit should be greatest precisely where accumulation is deepest (many
stacks) or data are scarcest (Milk, a single univariate series), which is
what the failure table shows. This picture is consistent with classical
observations on gradient-scale instability in deep networks; whether linear
output instability generalises beyond doubly residual signal routing
remains open.

\paragraph{Structural substitutes.} If that hypothesis is right, any
mechanism that inherently bounds output magnitude should make the
activation unnecessary --- and this is observed from two directions.
The AE trunk compresses to latent width $d$ before decoding; because
$d \ll U$, the decoded projection cannot produce arbitrarily large signals,
an implicit amplitude regulariser serving the same stabilising role through
structure (AE-family blocks on Milk show stuck rates in the low single
digits with the activation off, against 63--64\% for plain generic
counterparts). The frozen basis supplies the second instance: synthesis
through rows of unit $\ell_2$ norm lets the coefficient vector alone
control output energy, and the low-rank coefficient projection adds a
compression stage analogous to the bottleneck. Architectures with either
structural bound do not need the activation; the activation is best read as
an explicit, lightweight substitute for regularisation that bottlenecks and
orthonormal bases provide natively --- and its residual accuracy cost
(Table~\ref{tab:active_g_ablation}) says structural regularisation is the
more principled long-term solution.

\subsection{Tiered basis offsets: per-period verdict}
\label{app:ablations:tiered}

Tiering shifts each successive stack's window of the orthonormal basis
across the frequency-ordered spectrum instead of letting every stack read
the same leading rows. Table~\ref{tab:tiered_per_period} gives the
matched-seed comparison (paper-sample, plateau schedule) of tiered against
the closest non-tiered baseline per period. The verdict, consistent with
the main-text answer on when tiering could help: tiering does not help in
general. It survives only as a Daily-period intervention plus one isolated
significant Monthly cell; Weekly shows a genuine cascade reversal at its
horizon; and the earlier band-selection mechanism behind it holds only
under band-preserving constructions
(Table~\ref{tab:corrections_ledger}).

\begin{table}[h]
\centering
\small
\caption{Tiered basis offset, matched against the closest non-tiered
baseline. ``Pairs won'' counts wins for the tiered configuration out of 11
matched pairs; significance is paired Wilcoxon at $p < 0.05$.}
\label{tab:tiered_per_period}
\begin{tabular}{@{}llrll@{}}
\toprule
Period & Pairs won & Avg.\ $\Delta$ sMAPE & Sig.\ tiered $<$ base & Sig.\ base $<$ tiered \\
\midrule
Yearly    & $7/11$ & $-0.020$ & $1$ (alternating FB+T, $S{=}30$, AE) & $0$ \\
Quarterly & $8/11$ & $-0.020$ & $0$ & $0$ \\
Monthly   & $7/11$ & $-0.105$ & $1$ (alternating FB+T, $S{=}10$) & $0$ \\
Weekly    & $2/11$ & $\mathbf{+0.130}$ & $0$ & $0$ \\
Daily     & $\mathbf{11/11}$ & $\mathbf{-0.049}$ & $\mathbf{4}$ & $0$ \\
\bottomrule
\end{tabular}
\end{table}

The strongest single Daily wins (Cliff's $d \in [0.54, 0.78]$) are: the
gated-AE alternating pair at ten stacks ($\Delta = -0.080$, $p = 0.0036$,
$d = 0.78$); the AE-trunk alternating pair at thirty stacks
($\Delta = -0.074$, $p = 0.0036$, $d = 0.78$); a second gated-AE
alternating configuration at ten stacks ($\Delta = -0.067$,
$p = 0.0312$, $d = 0.58$); and the gated-AE alternating pair at thirty
stacks ($\Delta = -0.043$, $p = 0.0452$, $d = 0.54$). The Weekly reversal
was confirmed by an independent $n = 10$ plateau-schedule re-validation
released as \texttt{tiered\_offset\_m4\_weekly\_plateau\_validation\_results.csv}.

\subsection{LR scheduler: paper-sample matched pairs}
\label{app:ablations:lr}

Table~\ref{tab:lr_matched} gives the plateau-versus-MultiStepLR matched-pair
deltas on the 53-configuration paper-sample subset, the evidence behind the
per-period scheduler recommendations of Table~\ref{tab:lr_per_period}.
Negative deltas favor the plateau schedule.

\begin{table}[h]
\centering
\small
\caption{Paper-sample plateau-vs-step matched pairs
($\Delta =$ plateau $-$ step). Negative $\Delta$ favors plateau.}
\label{tab:lr_matched}
\begin{tabular}{@{}lrrr@{}}
\toprule
Period & Mean $\Delta$ & Median $\Delta$ & $n$ pairs \\
\midrule
Yearly    & $+0.007$ & $+0.021$ & $52$ \\
Quarterly & $\mathbf{-0.029}$ & $-0.032$ & $30$ \\
Monthly   & $\mathbf{-0.089}$ & $-0.198$ & $18$ \\
Weekly    & $+0.089$ & $+0.080$ & $11$ \\
Daily     & $\mathbf{-0.028}$ & $-0.028$ & $13$ \\
\bottomrule
\end{tabular}
\end{table}

The Hourly schedule choice is canonicalized at the cell level rather than
the period level: one tiered alternating configuration prefers plateau
(8.922 vs.\ 9.090, $p = 0.026$, $n = 10$), while a tiered AE-trunk unified
configuration prefers MultiStepLR (9.243 vs.\ 9.439, $p = 0.004$, $n = 10$).
The Hourly leaderboard cell therefore mixes scheduler choice by
configuration; this is why Table~\ref{tab:lr_per_period} records Hourly as
step-paper with the head-to-head not yet run across all novel variants.

\subsection{Divergence mechanism and rectifier diagnostic}
\label{app:ablations:divergence}

The divergence study behind Section~\ref{sec:results:stability} examined
what a diverged run actually computes. Across the forty-run diagnostic,
divergence is negative forecasts on positive series scored at sMAPE's 200
ceiling: deleting exactly those points recovers converged scores to within
0.11 sMAPE, so ``collapse'' is a scoring failure of a specific sign error,
not diffuse numerical blow-up. Rectifying the output head removes the
failure by construction on non-negative targets --- collapsed points fall
from 107 of 23{,}786 forecast points to 0 of 136 in the retrained arms ---
but rectification destroys the cancellation that sub-band basis designs
rely on, which is why it is not adopted globally and why the
$\dagger$-regime arms of the preprocessing comparison were discarded as an
unfair test (Section~\ref{sec:related:pramanick}). Per-run evidence is
released as \texttt{m4\_divergence\_signs\_results.csv} and
\texttt{m4\_dwt\_relu\_diagnostic\_results.csv}.

\subsection{Nonlinearity placement (boundary test)}
\label{app:ablations:boundary}

The completeness result of Section~\ref{sec:method_heads} requires linear
synthesis between the learned coefficient projection and the frozen basis:
with any interposed nonlinearity, the rotation-absorption argument no longer
applies. The boundary test checks the premise empirically by training
variants that insert an activation at different points along the synthesis
path --- after the trunk, between the coefficient projection and the basis,
and after the basis --- against the linear-synthesis reference.
Placing the nonlinearity between projection and basis measurably degrades
accuracy relative to the linear path, consistent with the scope condition;
placing it elsewhere does not engage the frozen-basis factorisation and
behaves like the corresponding unconstrained variant. We report the
qualitative ordering here and release the per-run record
(\texttt{m4\_p2\_boundary\_results.csv},
\texttt{m4\_dwt\_relu\_diagnostic\_results.csv}); no headline number in
this paper depends on the boundary arms.

\subsection{Traffic-96 NHiTS transfer detail}
\label{app:ablations:nhits}

Curated subsets of the block library were rerun on the NHiTS backbone for
Traffic at horizons $\{96, 192, 336, 720\}$ under MSE loss, 100 epochs, and
8 seeds, with pooling and hierarchical-interpolation schedules inherited
from the published recipe
(Section~\ref{sec:results:nhits}). Short-horizon ($H = 96$) winners are
unified frozen-basis-plus-trend blocks, where the lookback-to-horizon ratio
is largest and asymmetric backcast/forecast bases are most useful;
long-horizon ($H \geq 336$) winners are AE-trunk bottleneck variants. One
apparent winner involving a variational trunk is treated as NHiTS-specific
and carried into no recommendation. The scope caveat stands: a single
short-epoch regime at eight seeds carries high per-cell variance, so these
cells support only the interface-transfers claim, and the Traffic-horizon
detail stays here rather than in any leaderboard.

\subsection{What does not work}
\label{app:ablations:doesnotwork}

One row for completeness, since negative results are otherwise absent from
this paper. Pure variational trunks were catastrophic on M4 --- mean sMAPE
55--68 on Quarterly, Monthly, Weekly, and Daily --- across both the M4 and
Weather sweeps, and every variational row was removed from all leaderboards
after submission. We do not claim variational trunks are unworkable in
general, only that they were not stabilised within this study; the KL-based
basis construction that motivated them is developed separately in
Section~\ref{sec:method_basis}, where its theory-named optimum is evaluated
without variational training. The learned gate, by contrast, was shown null
rather than harmful (twelve comparisons, smallest $p = 0.24$), which is why
it is merged into the AE trunk category with a provenance marker instead of
dropped.


